\BourbakiExercisesSection

\begin{enumerate}
\def\labelenumi{\arabic{enumi}.}
\item
  Let E, F be two A-modules and \(u: \mathbf{E} \to \mathbf{F}\) a linear mapping. Let M be a submodule of E, N a submodule of F, \(i: \mathbf{M} \to \mathbf{E}\) the canonical injection and \(q: \mathbf{F} \to \mathbf{F} / \mathbf{N}\) the canonical surjection. Show that \(u(\mathbf{M}) = \operatorname{Im}(u \circ i)\) and \(u^{-1}(\mathbf{N}) = \operatorname{Ker}(q \circ u)\) .
\item
  Let E be a left A-module; for every left ideal a of A let aE denote the sum of the ax where \(x \in E\) , which is therefore a submodule of E.
\end{enumerate}

\begin{enumerate}
\def\labelenumi{(\alph{enumi})}
\item
  Show that for every family \((\mathbf{E}_{\lambda})_{\lambda \in L}\) of A-modules, a. \(\bigoplus_{\lambda \in L} \mathbf{E}_{\lambda}\) is canonically identified with \(\bigoplus_{\lambda \in L} a E_{\lambda}\) .
\item
  If \(\mathfrak{a}\) is a finitely generated left ideal of \(\mathbf{A}\) , show that for every family \((\mathbf{E}_{\lambda})_{\lambda \in L}\) of \(\mathbf{A}\) -modules, \(\mathfrak{a}\) . \(\prod_{\lambda \in L} \mathbf{E}_{\lambda}\) is canonically identified with \(\prod_{\lambda \in L} \mathfrak{a}\mathbf{E}_{\lambda}\) . Give an example where \(\mathfrak{a}\) is not finitely generated and the above property fails to hold (take \(\mathbf{A}\) commutative and all the \(\mathbf{E}_{\lambda}\) equal to \(\mathbf{A}\) ).
\end{enumerate}

*(c) Let K be a commutative field, A the ring K{[}X,Y{]} of polynomials in two indeterminates over K, m = AX and n = AY. If E is the quotient A-module of \(\mathbf{A}^2\) by the submodule of \(\mathbf{A}^2\) generated by \(\mathbf{Xe}_1 - \mathbf{Ye}_2\) (where \((e_1, e_2)\) is the canonical basis of \(\mathbf{A}^2\) ), show that \((m \cap n)E \neq (mE) \cap (nE)\) .*

\begin{enumerate}
\def\labelenumi{\arabic{enumi}.}
\setcounter{enumi}{2}
\tightlist
\item
  Let E be an A-module.
\end{enumerate}

\begin{enumerate}
\def\labelenumi{(\alph{enumi})}
\tightlist
\item
  Let \((\mathbf{L}_{\mu})_{\mu \in \mathbf{M}}\) be a family of right directed ordered sets and, for all \(\mu \in \mathbf{M}\) , let \((\mathrm{F}_{\lambda \mu})_{\lambda \in L_{\mu}}\) be an increasing family of submodules of \(\mathbf{E}\) . Show that
\end{enumerate}

\[
\bigcap_ {\mu \in \mathbf {M}} \left(\sum_ {\lambda \in L _ {\mu}} F _ {\lambda \mu}\right) = \sum_ {\rho \in \prod_ {\mu} L _ {\mu}} \left(\bigcap_ {\mu \in \mathbf {M}} F _ {\rho (\mu), \mu}\right).
\]

\begin{enumerate}
\def\labelenumi{(\alph{enumi})}
\setcounter{enumi}{1}
\tightlist
\item
  Take A to be a field and E the A-module \(A^{N}\) ; let G be the submodule \(\mathbf{A}^{(\mathbf{N})}\) of E; on the other hand, for every finite subset H of N, let \(F_{H}\) be the submodule \(\prod_{n\in \mathbf{N}}\mathbf{M}_n\) of \(\mathbf{E}\) , where \(\mathbf{M}_n = \mathbf{A}\) for \(n\notin \mathrm{H}\) and \(\mathbf{M}_n = 0\) for \(n\in \mathrm{H}\) . Show that the family \((\mathrm{F_H})\) , where \(\mathbf{H}\) runs throughout the set \(\Phi\) of finite subsets of \(\mathbf{N}\) is right directed and that
\end{enumerate}

\[
\left(\bigcap_ {\mathrm{H} \in \Phi} \mathrm{F} _ {\mathrm{H}}\right) + \mathrm{G} \neq \bigcap_ {\mathrm{H} \in \Phi} (\mathrm{F} _ {\mathrm{H}} + \mathrm{G}).
\]

\begin{enumerate}
\def\labelenumi{\arabic{enumi}.}
\setcounter{enumi}{3}
\tightlist
\item
  Let \(\mathbf{E}, \mathbf{F}_1, \mathbf{F}_2\) be three A-modules and \(f_1: \mathbf{F}_1 \to \mathbf{E}\) , \(f_2: \mathbf{F}_2 \to \mathbf{E}\) two A-linear mappings. The fibre product of \(\mathbf{F}_1\) and \(\mathbf{F}_2\) relative to \(f_1\) and \(f_2\) is the submodule of the product \(\mathbf{F}_1 \times \mathbf{F}_2\) consisting of the ordered pairs \((x_1, x_2)\) such that \(f_1(x_1) = f_2(x_2)\) ; it is denoted by \(\mathbf{F}_1 \times_{\mathbf{E}} \mathbf{F}_2\) .
\end{enumerate}

\begin{enumerate}
\def\labelenumi{(\alph{enumi})}
\item
  For every ordered pair of A-linear mappings \(u_1 \colon G \to F_1, u_2 \colon G \to F_2\) such that \(f_1 \circ u_1 = f_2 \circ u_2\) , there exists one and only one A-linear mapping \(u \colon G \to F_1 \times_E F_2\) such that \(p_1 \circ u = u_1, p_2 \circ u = u_2\) , where \(p_1\) and \(p_2\) denote the restrictions to \(F_1 \times_E F_2\) of the projections \(pr_1\) and \(pr_2\) .
\item
  Let \(\mathbf{E}'\) , \(\mathbf{F}_1'\) , \(\mathbf{F}_2'\) be three A-modules, \(f_1': \mathbf{F}_1' \to \mathbf{E}', f_2': \mathbf{F}_2' \to \mathbf{E}'\) A-linear mappings and \(\mathbf{F}_1' \times_{\mathbf{E}'} \mathbf{F}_2'\) the fibre product relative to these mappings. For every system of A-linear mappings \(v_1: \mathbf{F}_1 \to \mathbf{F}_1'\) , \(v_2: \mathbf{F}_2 \to \mathbf{F}_2'\) , \(w: \mathbf{E} \to \mathbf{E}'\) such that \(f_1' \circ v_1 = w \circ f_1, f_2' \circ v_2 = w \circ f_2\) , let \(v\) be the A-linear mapping of \(\mathbf{F}_1 \times_{\mathbf{E}} \mathbf{F}_2\) into \(\mathbf{F}_1' \times_{\mathbf{E}'} \mathbf{F}_2'\) corresponding to the two linear mappings \(v_1 \circ p_1\) and \(v_2 \circ p_2\) . Show that if \(v_1\) and \(v_2\) are injective, so is \(v\) ; give an example where \(v_1\) and \(v_2\) are surjective but \(v\) is not surjective (take \(\mathbf{E}' = \{0\}\) ). If \(w\) is injective and \(v_1\) and \(v_2\) surjective, show that \(v\) is surjective.
\end{enumerate}

\begin{enumerate}
\def\labelenumi{\arabic{enumi}.}
\setcounter{enumi}{4}
\tightlist
\item
  Let \(\mathbf{E}, \mathbf{F}_1, \mathbf{F}_2\) be three A-modules and \(f_1: \mathbf{E} \to \mathbf{F}_1, f_2: \mathbf{E} \to \mathbf{F}_2\) two linear mappings. The amalgamated sum of \(\mathbf{F}_1\) and \(\mathbf{F}_2\) relative to \(f_1\) and \(f_2\) is the quotient module of \(\mathbf{F}_1 \times \mathbf{F}_2\) by the submodule the image of \(\mathbf{E}\) under the mapping \(z \mapsto (f_1(z), -f_2(z))\) ; it is denoted by \(\mathbf{F}_1 \oplus_{\mathbf{E}} \mathbf{F}_2\) .
\end{enumerate}

\begin{enumerate}
\def\labelenumi{(\alph{enumi})}
\tightlist
\item
  For every ordered pair of A-linear mappings \(u_{1}: \mathbf{F}_{1} \to \mathbf{G}, u_{2}: \mathbf{F}_{2} \to \mathbf{G}\) such that \(u_{1} \circ f_{1} = u_{2} \circ f_{2}\) , show that there exists one and only one A-linear mapping \(u: \mathbf{F}_{1} \oplus_{\mathbf{E}} \mathbf{F}_{2} \to \mathbf{G}\) such that \(u \circ j_{1} = u_{1}, u \circ j_{2} = u_{2}\) , where \(j_{1}\) and \(j_{2}\) denote the respective compositions of the canonical mapping
\end{enumerate}

\[
\mathbf {F} _ {1} \times \mathbf {F} _ {2} \rightarrow \mathbf {F} _ {1} \oplus_ {\mathrm{E}} \mathbf {F} _ {2}
\]

with the canonical injections \(\mathbf{F}_1\to \mathbf{F}_1\times \mathbf{F}_2,\mathbf{F}_2\rightarrow \mathbf{F}_1\times \mathbf{F}_2.\)

\begin{enumerate}
\def\labelenumi{(\alph{enumi})}
\setcounter{enumi}{1}
\tightlist
\item
  Let \(\mathbf{E}'\) , \(\mathbf{F}_1'\) , \(\mathbf{F}_2'\) be three A-modules, \(f_1': \mathbf{E}' \to \mathbf{F}_1'\) , \(f_2': \mathbf{E}' \to \mathbf{F}_2'\) two A-linear mappings and \(\mathbf{F}_1' \oplus_{\mathbf{E}'} \mathbf{F}_2'\) the amalgamated sum relative to these mappings. For every system of A-linear mappings \(v_1: \mathbf{F}_1' \to \mathbf{F}_1\) , \(v_2: \mathbf{F}_2' \to \mathbf{F}_2\) , \(w: \mathbf{E}' \to \mathbf{E}\) such that \(v_1 \circ f_1' = f_1 \circ w\) , \(v_2 \circ f_2' = f_2 \circ w\) , let \(v\) be the A-linear mapping of \(\mathbf{F}_1' \oplus_{\mathbf{E}'} \mathbf{F}_2'\) into \(\mathbf{F}_1 \oplus_{\mathbf{E}} \mathbf{F}_2\) corresponding to the two linear mappings \(j_1 \circ v_1\) and \(j_2 \circ v_2\) . Show that if \(v_1\) and \(v_2\) are surjective, so is \(v\) ; give an example where \(v_1\) and \(v_2\) are injective but \(v\) is not injective. Show that if \(w\) is surjective and \(v_1\) and \(v_2\) injective, then \(v\) is injective.
\end{enumerate}

\begin{enumerate}
\def\labelenumi{\arabic{enumi}.}
\setcounter{enumi}{5}
\tightlist
\item
  Given an A-module E, let \(\gamma(E)\) denote the least of the cardinals of the generating systems of E.
\end{enumerate}

\begin{enumerate}
\def\labelenumi{(\alph{enumi})}
\tightlist
\item
  If \(0 \to \mathbf{E} \to \mathbf{F} \to \mathbf{G} \to 0\) is an exact sequence of A-modules, then
\end{enumerate}

\[
\gamma (\mathbf {G}) \leqslant \gamma (\mathbf {F}) \leqslant \gamma (\mathbf {E}) + \gamma (\mathbf {G}).
\]

Give an example of a product \(\mathbf{F} = \mathbf{E} \times \mathbf{G}\) such that \(\gamma(\mathbf{E}) = \gamma(\mathbf{F}) = \gamma(\mathbf{G}) = 1\) (take \(\mathbf{E}\) and \(\mathbf{G}\) to be quotients of \(\mathbf{Z}\) ).

\begin{enumerate}
\def\labelenumi{(\alph{enumi})}
\setcounter{enumi}{1}
\item
  If \(\mathbf{E}\) is the sum of a family \((\mathrm{F}_{\lambda})_{\lambda \in L}\) of submodules, then \(\gamma (\mathbf{E})\leqslant \sum_{\lambda \in L}\gamma (\mathbf{F}_{\lambda}).\)
\item
  If \(\mathbf{M}\) and \(\mathbf{N}\) are submodules of a module \(\mathbf{E}\) , show that
\end{enumerate}

\[
\sup (\gamma (\mathbf {M}), \gamma (\mathbf {N})) \leqslant \gamma (\mathbf {M} \cap \mathbf {N}) + \gamma (\mathbf {M} + \mathbf {N}).
\]

\begin{enumerate}
\def\labelenumi{\arabic{enumi}.}
\setcounter{enumi}{6}
\tightlist
\item
  \begin{enumerate}
  \def\labelenumii{(\alph{enumii})}
  \tightlist
  \item
    Let A be a ring and E an (A, A)-bimodule; on the product \(B = A \times E\) a ring structure is defined by setting
  \end{enumerate}
\end{enumerate}

\[
(a, x) (a ^ {\prime}, x ^ {\prime}) = (a a ^ {\prime}, a x ^ {\prime} + x a ^ {\prime});
\]

E (canonically identified with \(\{0\} \times \mathbf{E}\) ) is then a two-sided ideal of B such that \(\mathbf{E}^2 = \{0\}\) .

\begin{enumerate}
\def\labelenumi{(\alph{enumi})}
\setcounter{enumi}{1}
\tightlist
\item
  By suitably choosing A and E, show that E can be a left B-module with \(\gamma(\mathbf{E})\) (Exercise 6) an arbitrary cardinal, although \(\gamma(\mathbf{B}_{s}) = 1\) .
\end{enumerate}

*8. Let A be a commutative ring, \(\mathbf{B} = \mathbf{A}[\mathbf{X}_n]_{n\geqslant 1}\) a polynomial ring over A in an infinity of indeterminates and C the quotient ring of B by the ideal generated by the polynomials \((\mathbf{X}_1 - \mathbf{X}_2)\mathbf{X}_i\) for \(i\geqslant 3\) . If \(\xi_1\) and \(\xi_2\) are the classes in C of \(\mathbf{X}_1\) and \(\mathbf{X}_2\) respectively, show that the intersection in C of the monogenous ideals \(\mathrm{C}\xi_1\) and \(\mathrm{C}\xi_2\) is not a finitely generated ideal.*

\begin{enumerate}
\def\labelenumi{\arabic{enumi}.}
\setcounter{enumi}{8}
\item
  Show that in an A-module E, if \((\mathrm{F}_{n})_{n\geqslant1}\) is a strictly increasing sequence of finitely generated submodules, the submodule F of E the union of the \(F_{n}\) is not finitely generated. Deduce that if E is an A-module which is not finitely generated, there exists a submodule F of E such that \(\gamma(\mathrm{F})=\aleph_{0}(=\mathrm{Card}(\mathbf{N}))\) .
\item
  Let E be a Z-module the direct sum of two submodules M, N respectively isomorphic to Z/2Z and Z/3Z; show that there exists no other decomposition of E as a direct sum of two non-zero submodules.
\item
  Let A be a ring admitting a field of left fractions (I, § 9, Exercise 15). Show that in the A-module \(A_{s}\) , a submodule distinct from \(A_{s}\) and \(\{0\}\) admits no supplementary submodule (cf.~I, § 9, Exercise 16).
\item
  Let \(\mathbf{G}\) be a module and \(\mathbf{E}, \mathbf{F}\) two submodules such that \(\mathbf{E} \subset \mathbf{F}\) .
\end{enumerate}

\begin{enumerate}
\def\labelenumi{(\alph{enumi})}
\item
  If F is a direct factor of G, F/E is a direct factor of G/E. If also E is a direct factor of F, E is a direct factor of G.
\item
  If E is a direct factor of G, then E is a direct factor of F. If also F/E is a direct factor of G/E, then F is a direct factor in G.
\item
  Give an example of two submodules \(\mathbf{M}, \mathbf{N}\) of the \(\mathbf{Z}\) -module \(\mathbf{E} = \mathbf{Z}^2\) such that \(\mathbf{M}\) and \(\mathbf{N}\) are direct factors of \(\mathbf{E}\) but \(\mathbf{M} + \mathbf{N}\) is not a direct factor of \(\mathbf{E}\) .
\item
  Let \(p\) be a prime number and \(U_p\) the submodule of the \(\mathbf{Z}\) -module \(\mathbf{Q} / \mathbf{Z}\) consisting of the classes mod \(\mathbf{Z}\) of rational numbers of the form \(k / p^n\) ( \(k \in \mathbf{Z}\) , \(n \in \mathbf{N}\) ). Let \(E\) be the product \(\mathbf{Z}\) -module \(M \times N\) , where \(M\) and \(N\) are both isomorphic to \(U_p\) , \(M\) and \(N\) being canonically identified with submodules of \(E\) . Consider the endomorphism \(u: (x, y) \mapsto (x, y + px)\) of \(E\) ; show that \(u\) is bijective. If \(M' = u(M)\) , show that \(M\) and \(M'\) are both direct factors of \(E\) but that \(M \cap M'\) is not a direct factor of \(E\) (use (b), showing that the \(Z\) -module \(U_p\) has no direct factor other than itself and \{0\}).
\end{enumerate}

¶ 13. Let E be an A-module and B the ring End\_A E. Show that for an element \(u \in B\) the following three conditions are equivalent: (1) Bu is a direct factor of the left B-module B\_s; (2) uB is a direct factor of the right B-module B\_d; (3) Ker(u) and Im(u) are direct factors of the A-module E. (Observe that every left ideal which is a direct factor of B\_s is of the form Bp, where p is a projector in E; cf.~I, § 8, Exercise 12 and use II, § 1, no. 9, Corollaries 1 and 2 to Proposition 15.)

\begin{enumerate}
\def\labelenumi{\arabic{enumi}.}
\setcounter{enumi}{13}
\tightlist
\item
  Let L be a well-ordered set, E an A-module and \((\mathbf{E}_{\lambda})_{\lambda \in L}\) an increasing family of submodules of E such that: (1) there exists a \(\lambda \in L\) such that \(E_{\lambda} = \{0\}\) ; (2) E is the union of the \(E_{\lambda}\) for \(\lambda \in L\) ; (3) if \(\lambda \in L\) is such that the set of \(\mu < \lambda\) admits a greatest element \(\lambda'\) , \(E_{\lambda}\) is the direct sum \(E_{\lambda'}\) and a submodule \(F_{\lambda}\) ; (4) if \(\lambda \in L\) is the least upper bound of the set of \(\mu < \lambda\) , then \(E_{\lambda}\) is the union of the \(E_{\mu}\) for \(\mu < \lambda\) . Under these conditions, show that E is the direct sum of the \(F_{\lambda}\) . (Prove by transfinite induction that each \(E_{\lambda}\) is the direct sum of the \(F_{\mu}\) such that \(\mu \leqslant \lambda\) .)
\end{enumerate}

¶ 15. Let c be an infinite cardinal.

\begin{enumerate}
\def\labelenumi{(\alph{enumi})}
\item
  Let I be a set and \(R\{a, b\}\) a relation between elements of I such that for all \(\alpha \in I\) the set of \(\beta \in I\) for which \(R\{\alpha, \beta\}\) holds has cardinal \(\leqslant c\) . Show that there exists a well-ordered set L and an increasing family \((\mathrm{I}_{\lambda})_{\lambda \in L}\) of subsets of I such that: (1) there exists a \(\lambda \in L\) such that \(I_{\lambda} = \varnothing\) ; (2) I is the union of the \(I_{\lambda}\) for \(\lambda \in L\) ; (3) if \(\lambda \in L\) is such that the set of \(\mu < \lambda\) has a greatest element \(\lambda'\) , \(\text{Card}(\mathbf{I}_{\lambda} - \mathbf{I}_{\lambda'}) \leqslant c\) ; (4) if \(\lambda \in L\) is the least upper bound of the set of \(\mu < \lambda\) , \(I_{\lambda}\) is the union of the \(I_{\mu}\) for \(\mu < \lambda\) ; (5) if \(\alpha \in I_{\lambda}\) and \(R\{\alpha, \beta\}\) holds, then \(\beta \in I_{\lambda}\) . (Note first that for all \(\alpha \in I\) the set of \(\beta \in I\) for which there exists an integer n and a sequence \((\gamma_{j})_{1 \leqslant j \leqslant n}\) of elements of I such that \(\gamma_{1} = \alpha, \gamma_{n} = \beta\) and \(R\{\gamma_{j}, \gamma_{j+1}\}\) holds for \(1 \leqslant j \leqslant n - 1\) , has cardinal \(\leqslant c\) . Then construct the \(I_{\lambda}\) by transfinite induction.)
\item
  Let E be an A-module the direct sum of a family \((\mathbf{M}_{\alpha})_{\alpha\in I}\) of submodules such that \(\gamma(\mathbf{M}_{\alpha})\leqslant c\) for all \(\alpha\in I\) (cf.~Exercise 6). Let f be an endomorphism of E. Show that there exists a well-ordered set L and an increasing family \((\mathbf{I}_{\lambda})_{\lambda\in L}\) of subsets of I satisfying properties (1), (2), (3) and (4) of (a) and such moreover that if we write \(\mathbf{E}_{\lambda} = \bigoplus_{\alpha \in \mathbf{I}_{\lambda}} \mathbf{M}_{\alpha}\) , then \(f(\mathbf{E}_{\lambda}) \subset \mathbf{E}_{\lambda}\) . (Apply (a) to the relation \(R\{\alpha, \beta\}\) : ``there exists \(x \in \mathbf{M}_{\alpha}\) such that the component of \(f(x)\) in \(\mathbf{M}_{\beta}\) is \(\neq 0\)''.) For all \(\lambda \in \mathbf{L}\) such that the set of \(\mu < \lambda\) has a greatest element \(\lambda'\) , write \(\mathbf{F}_{\lambda} = \bigoplus_{\alpha \in \mathbf{I}_{\lambda} - \mathbf{I}_{\lambda'}} \mathbf{M}_{\alpha}\) . Show that \(\mathbf{E}\) is the direct sum of the \(\mathbf{F}_{\lambda}\) (apply Exercise 14).
\item
  With the hypotheses and notation of (b), suppose further that \(f\) is a projector and write \(\mathbf{P} = f(\mathbf{E})\) . Let \(\mathbf{P}_{\lambda} = \mathbf{P} \cap \mathbf{E}_{\lambda}\) ; show that if \(\lambda \in \mathbf{L}\) is such that the set of \(\mu < \lambda\) has a greatest element \(\lambda'\) , \(\mathbf{P}_{\lambda'}\) is a direct factor of \(\mathbf{P}_{\lambda}\) (cf.~Exercise 12), that a supplementary submodule \(\mathbf{P}_{\lambda}'\) of \(\mathbf{P}_{\lambda'}\) in \(\mathbf{P}_{\lambda}\) is isomorphic to a direct factor of \(\mathbf{F}_{\lambda}\) and that \(\gamma(\mathbf{P}_{\lambda}') \leqslant c\) . Show that \(\mathbf{P}\) is the direct sum of the \(\mathbf{P}_{\lambda}'\) (apply Exercise 14).
\item
  Deduce from (c) that if an A-module E is the direct sum of a family \((\mathrm{M}_{\alpha})_{\alpha\in I}\) of submodules such that \(\gamma(\mathrm{M}_{\alpha})\leqslant\mathfrak{c}\) , then every direct factor P of E is also the direct sum of a family \((\mathrm{N}_{\lambda})_{\lambda\in L}\) of submodules such that \(\gamma(\mathrm{N}_{\lambda})\leqslant\mathfrak{c}\) .
\end{enumerate}

\begin{enumerate}
\def\labelenumi{\arabic{enumi}.}
\setcounter{enumi}{15}
\tightlist
\item
  Let A be a ring such that there exists an A-module M with a generating system of n elements but containing a free system of \(n + 1\) elements.
\end{enumerate}

\begin{enumerate}
\def\labelenumi{(\alph{enumi})}
\item
  Show that there exists in \(\mathbf{A}_s^n\) a free system of \(n + 1\) elements.
\item
  Deduce that there exists in M an infinite free system (construct such a system by induction using (a)).
\item
  Let C be a ring, E the free C-module \(\mathrm{C}_{s}^{(\mathrm{N})}, (e_{n})\) its canonical basis and A the endomorphism ring of E. Let \(u_{1}\) and \(u_{2}\) denote the endomorphisms of E defined by the conditions \(u_{1}(e_{2n}) = e_{n}, u_{1}(e_{2n+1}) = 0, u_{2}(e_{2n+1}) = e_{n}, u_{2}(e_{2n}) = 0\) for all \(n \geqslant 0\) . Show that \(u_{1}\) and \(u_{2}\) form a basis of the A-module \(A_{s}\) and deduce that in the A-module \(A_{s}\) there exist infinite free systems.
\end{enumerate}

*(d) Let A be the tensor algebra of a vector space E of dimension \(\geqslant 2\) over a commutative field (cf.~III, § 5, no. 1). Show that in the A-module \(A_s\) there exist infinite free systems (observe that two linearly independent vectors in E are linearly independent in \(A_s\) and use \((b)\) ). On the other hand, for every integer \(n > 0\) , every basis of the A-module \(A_s^n\) has \(n\) elements (cf.~§ 7, no. 2, Proposition 3).*

\begin{enumerate}
\def\labelenumi{\arabic{enumi}.}
\setcounter{enumi}{16}
\tightlist
\item
  Let A be a ring.
\end{enumerate}

\begin{enumerate}
\def\labelenumi{(\alph{enumi})}
\item
  Show that if the A-module \(\mathbf{A}_s^n\) is isomorphic to no \(\mathbf{A}_s^m\) for \(m > n\) , then, for all \(p < n\) , \(\mathbf{A}_s^p\) is isomorphic to no \(\mathbf{A}_s^q\) for \(q > p\) .
\item
  Show that if the A-module \(A_{s}^{n}\) contains no free system of \(n + 1\) elements, the A-module \(A_{s}^{p}\) contains no free system of \(p + 1\) elements for p \textless{} n.
\end{enumerate}

\begin{enumerate}
\def\labelenumi{\arabic{enumi}.}
\setcounter{enumi}{17}
\tightlist
\item
  Let A be a ring for which there exists an integer \(p\) with the following property: for every family \((a_i)_{1 \leqslant i \leqslant p}\) of \(p\) elements of A, there exists a family \((c_i)_{1 \leqslant i \leqslant p}\) of elements of A, one at least of which is not a right divisor of 0, such that \(\sum_{i=1}^{p} c_i a_i = 0\) .
\end{enumerate}

\begin{enumerate}
\def\labelenumi{(\alph{enumi})}
\item
  Show, by induction on \(n\) , that for every family \((x_{i})\) of \(p^n\) elements of \(A_s^n\) , there exists a family \((c_j)\) of \(p^n\) elements of \(A\) , one at least of which is not a right divisor of 0, such that \(\sum_{j} c_j x_j = 0\) .
\item
  Deduce from (a) and Exercise 16 that for all n \textgreater{} 0, the A-module \(A_{s}^{n}\) contains no free system of \(n + 1\) elements.
\end{enumerate}

\begin{enumerate}
\def\labelenumi{\arabic{enumi}.}
\setcounter{enumi}{18}
\tightlist
\item
  Let A be a non-zero ring with no divisors of zero.
\end{enumerate}

\begin{enumerate}
\def\labelenumi{(\alph{enumi})}
\item
  Show that for \(n > 1\) , the A-module \(\mathbf{A}_s^n\) is not monogenous.
\item
  Given an integer \(n \geqslant 1\) , for the A-module \(\mathbf{A}_s^n\) to contain no free system of \(n + 1\) elements, it is necessary that A admit a field of left fractions (I, §9, Exercise 15); conversely, if this condition is satisfied, \(\mathbf{A}_s^n\) contains a free system of \(n + 1\) elements for no integer \(n > 0\) . (Use Exercises 17 and 18 and I, §9, Exercise 16.)
\item
  Show that if every left ideal of A is monogenous, A admits a field of left fractions (use (a) and I, § 9, Exercise 16).
\end{enumerate}

\begin{enumerate}
\def\labelenumi{\arabic{enumi}.}
\setcounter{enumi}{19}
\tightlist
\item
  \begin{enumerate}
  \def\labelenumii{(\alph{enumii})}
  \tightlist
  \item
    Let A be a ring admitting a field of left fractions (I, § 9, Exercise 15). Let E be an A-module; show that, if in E \((x_{i})_{1\leqslant i\leqslant r}\) is a free system and \(y,z\) two elements such that \(x_{1},\ldots ,x_{r},y\) on the one hand and \(x_{1},\ldots ,x_{r},z\) on the other are two related systems, then \(x_{2},\ldots ,x_{r},y,z\) is a related system.
  \end{enumerate}
\end{enumerate}

\begin{enumerate}
\def\labelenumi{(\alph{enumi})}
\setcounter{enumi}{1}
\tightlist
\item
  Let A be the quotient ring Z/6Z and let \((e_{1}, e_{2})\) be the canonical basis of \(A^{2}\) ; show that if \(a = 2e_{1} + 3e_{2}\) , a and \(e_{1}\) form a related system, as do a and \(e_{2}\) , although \(e_{1}\) and \(e_{2}\) form a free system.
\end{enumerate}

\begin{enumerate}
\def\labelenumi{\arabic{enumi}.}
\setcounter{enumi}{20}
\item
  Let A be a ring, b, c two elements of A which are not right divisors of 0, M the A-module A/Abc and N the submodule Ac/Abc of M. For N to be a direct factor in M, it is necessary and sufficient that there exist two elements x, y of A such that \(xb + cy = 1\) . (If \(\varepsilon\) is the class of 1 in M, show that a supplementary subspace of N in M is generated by an element of the form \((1 - yc)\varepsilon\) whose annihilator is the ideal Ac.)
\item
  If an A-module E admits a basis whose indexing set is I and one of the two sets A, I is infinite, show that E is equipotent to \(A \times I\) .
\item
  \begin{enumerate}
  \def\labelenumii{(\alph{enumii})}
  \tightlist
  \item
    Let M, N be two subsets of an A-module E and m and n their annihilators; show that the annihilator of M ∩ N contains m + n and give an example where it is distinct from m + n.
  \end{enumerate}
\end{enumerate}

\begin{enumerate}
\def\labelenumi{(\alph{enumi})}
\setcounter{enumi}{1}
\item
  In a product module \(\prod_{i} \mathbf{E}_{i}\) , the annihilator of a subset \(\mathbf{F}\) is the intersection of the annihilators of its projections.
\item
  In a free A-module, the annihilator of an element \(\neq0\) contains only left divisors of 0 in A; in particular, if A is a ring with no divisors of 0, every element \(\neq0\) of a free module is free.
\end{enumerate}

\begin{enumerate}
\def\labelenumi{\arabic{enumi}.}
\setcounter{enumi}{23}
\tightlist
\item
  Let E be an A-module and M, N two submodules of E. The transporter of M into N, denoted by N:M, is the set of \(a \in A\) such that \(aM \subset N\) ; it is a two-sided ideal of A, equal to \(\text{Ann}((\mathbf{M} + \mathbf{N})/\mathbf{N})\) .
\end{enumerate}

\begin{enumerate}
\def\labelenumi{(\alph{enumi})}
\item
  Let A be a commutative ring, M an A-module, N a sub-module of M and x an element of M; show that \(N \cap Ax = (N:Ax)x\) .
\item
  Let A be a commutative ring and \(a, b\) two ideals of A. Show that \(a:b \supset a\) and that \((a:b)/a\) is an A-module isomorphic to \(\operatorname{Hom}(A/b, A/a)\) .
\end{enumerate}

\begin{enumerate}
\def\labelenumi{\arabic{enumi}.}
\setcounter{enumi}{24}
\item
  Let E, F be two A-modules and \(u: E \to F\) a linear mapping. Show that the mapping \((x, y) \mapsto (x, y - u(x))\) of the product module \(E \times F\) to itself is an automorphism of \(E \times F\) . Deduce that if there exists a linear mapping \(v: F \to E\) and an \(a \in E\) such that \(v(u(a)) = a\) , there exists an automorphism w of \(E \times F\) such that \(w(a, 0) = (0, u(a))\) .
\item
  Let E be a free A-module with a basis containing at least two elements.
\end{enumerate}

\begin{enumerate}
\def\labelenumi{(\alph{enumi})}
\item
  Show that every semi-linear mapping of E into itself (relative to an automorphism of A) which permutes with all the automorphisms of E, is necessarily a homothety \(x \mapsto \alpha x (\alpha \in \mathbf{A})\) .
\item
  Deduce from (a) that the centre of \(\operatorname{End}_{\mathbf{A}}(\mathbf{E})\) is the ring of central homotheties, isomorphic to the centre of A.
\item
  Deduce from (a) that the centre of the group \(\mathbf{GL}(\mathbf{E})\) is the group of invertible central homotheties of E, which is isomorphic to the multiplicative group of invertible elements of the centre of A.
\end{enumerate}

\begin{enumerate}
\def\labelenumi{\arabic{enumi}.}
\setcounter{enumi}{26}
\tightlist
\item
  Let A be a commutative ring. Show that if an ideal \(\mathfrak{J}\) of A is a free A-module, then \(\mathfrak{J}\) is a monogenous A-module (in other words a principal ideal). Give an example showing that the proposition does not extend to the left ideals in a non-commutative ring (Exercise 16).
\end{enumerate}

